\documentclass[10pt]{article}
\usepackage[margin=0.8in]{geometry}
\usepackage{amsmath,microtype}
\begin{document}
\begin{center}
{\LARGE Guide: Merging Paths of Different Lengths}\par\medskip
Collatz proof project\quad---\quad September 5, 2026
\end{center}
\raggedright
\section*{What changed}
Earlier replacements required the two paths to take the same number of
steps. Allowing different lengths finds additional smaller starting seeds
that reach the same endpoint. Lean now checks when such a replacement
preserves noncontraction for the entire future orbit.

\section*{What noncontraction means}
After $t$ shortcut steps with $j$ odd steps, the multiplicative coefficient
is $3^j/2^t$. A seed is called noncontracting when that coefficient is at
least one for every prefix. This is stronger than merely seeing a finite
initial run without descent. No infinite noncontracting seed is constructed.

\section*{The two checks needed for transport}
The smaller seed's path must have coefficient at least one at every step
before the meeting. At the meeting, its coefficient must be at least as
large as the original path's coefficient. Afterward the two paths follow
the same values, so they receive the same further multiplicative factors.
Those two checks preserve full noncontraction if the original seed has it.

This matters because shifting a noncontracting seed forward does not
automatically give another noncontracting seed. The comparison at the
meeting accounts for the coefficient already accumulated before that shift.

\section*{A concrete infinite family}
For any whole number $Q\geq0$, take
\[
N=111+4374Q,\qquad M=103+4096Q.
\]
The larger seed reaches $167+6561Q$ in one step. The smaller seed reaches
the same value in twelve steps with eight odd steps. Its prefix passes
both transport checks. For $Q=0$, this is the meeting of 111 and 103 at 167.
The difference $N-M=8+278Q$ is always positive.

Thus, if $N$ were noncontracting, the smaller $M$ would be noncontracting
too. A least positive noncontracting seed cannot be in this progression.
This does not show that every member converges, or cover all possible
least seeds. The meeting occurs before any forward descent of $N$.

\section*{What the search establishes}
Among 7,495 seeds below $2^{18}$ with noncontracting prefixes through time
18, the previous equal-time/equal-count test covers 1,391 as smaller
merges. Searching different lengths through 18 adds 3,768, all with
certificates satisfying the new transport theorem. There are still 2,336
unresolved seeds within this search, including 27 and 703.

These counts concern finitely many starting seeds and bounded path
lengths. They are exact Python results, not kernel-certified global
coverage. An unresolved seed can have a useful longer replacement;
it is not a counterexample to Collatz.

\section*{Reproduce and continue}
From \texttt{lean/}, run \texttt{lake build Collatz.DominatingMerge}.
The search is \texttt{analysis/variable\_length\_surgery.py}, with a
matching test file and saved results. The next challenge is to find
arithmetic conditions that guarantee a suitable replacement beyond these
bounded searches. Neither a proof of nondivergence nor a full Collatz
proof has been obtained.
\end{document}
